\section{Introduction}

Two-way deterministic automata recognize the same languages as one-way
finite automata, as established by Rabin and Scott and by
Shepherdson~\cite{RabinScott1959,Shepherdson1959}. Allowing the head to
return to a symbol can nevertheless change the amount of finite control
needed to recognize a language. How much finite control is needed to
replace nondeterministic choices by deterministic head motion?
Sakoda and Sipser posed the state cost of converting one-way and two-way
nondeterministic finite automata to two-way deterministic automata in
1978~\cite{SakodaSipser1978}.
The question concerns succinctness: the simulator may revisit the input,
but must store all information other than the head position in its finite
state. A polynomial simulation would make nondeterminism inexpensive in
this measure even when the source machine itself moves in both directions.

The witness used here is Sakoda and Sipser's complete family $B_h$,
now called \emph{one-way liveness}: each letter specifies a binary
relation on an $h$-element set, and a word is live when the product of
its letters is nonempty.
Their Theorem~2.3 establishes this family's completeness for
one-way-nondeterministic to two-way-deterministic state
conversion~\cite{SakodaSipser1978}. Our proof below gives the small
source automaton and the context-separation property directly, with
the conventions and exact state counts used here.

Earlier exponential lower bounds restrict the motion of the deterministic
target. Sipser proved an exponential separation for sweeping automata,
whose head changes direction only at an endmarker~\cite{Sipser1980}.
Kapoutsis extended this phenomenon to deterministic machines constrained
to make sublinearly many reversals as a function of input
length~\cite{Kapoutsis2013}. An unrestricted
target can instead reverse at an interior cell and make arbitrarily many
reversals as the input length grows.

Quadratic lower bounds against unrestricted two-way deterministic targets
already occur in Chrobak's work on unary
languages~\cite{Chrobak1986Unary,Chrobak2003Errata}. For the promise restriction of
one-way liveness to words of exactly three relation letters, Kapoutsis
proved the optimal order $\Theta(h^2/\log h)$~\cite{Kapoutsis2018ShortLiveness};
the alphabet there is still the full relation alphabet. Adeogun and
Kapoutsis~\cite{AdeogunKapoutsis2026} prove a quadratic lower bound for
unrestricted deterministic two-way automata against one-way liveness.
We prove an exponential bound for that unrestricted target model.
The source automata use no left moves, whereas the lower bound permits
stay moves, partial transition rules, and nonaccepting infinite
computations in the target.

The main representation difficulty is to account for repeated visits to
a cell while assigning each input symbol a fixed object independent of
its neighbors. We use the Brauer diagram monoid, whose matchings
underlie Brauer's algebras~\cite{Brauer1937,Auinger2012}:
its elements are perfect matchings on two equally sized sets of ports,
and multiplication glues the adjoining sets and contracts paths.
Determinism makes the configurations reaching a designated accepting
sink form a tree. A tour of that tree gives the required matching
representation, with separate leaves for unused transition candidates
ensuring that each cell's wiring depends only on its own symbol.
The algebraic part then proves an exponential lower bound on the number
of ports whenever matching products retain all the information in
arbitrary relation products. All diagram and relation-monoid facts
needed for this argument are proved below.

\subsection{Model and statement}

An automaton has a finite state set and one read-only head on a finite word
over an alphabet $\Sigma$, bracketed by two distinct endmarkers outside
$\Sigma$. Its head starts at the left endmarker in a designated initial
state. A transition changes the state and moves the head left, right, or
not at all; the head cannot cross an endmarker. A transition depends only
on the current state and scanned symbol. Deterministic transition rules
may be partial. A nondeterministic rule permits several successors.
There are no additional heads, writable cells, work tapes, advice, or
randomness. All states, including initial and accepting states, are counted.

Acceptance means that a finite computation reaches an accepting state;
a nonaccepting infinite computation rejects. We initially count an
accepting initial configuration as success. If success instead requires
a positive transition into an accepting state, the normalization in
Section~\ref{sec:machines} uses one additional initial copy. The main
bound below includes this extra state and holds with either convention.

For a finite set $H$, let $\R{H}$ be the monoid of all binary relations
on $H$. Products are in path order from left to right:
\[
 (x,z)\in AB
 \quad\Longleftrightarrow\quad
 \text{there exists }y\in H\text{ with }(x,y)\in A,\ (y,z)\in B.
\]
Write $\Id{F}=\{(x,x):x\in F\}$ for $F\subseteq H$; the identity is
$\Id{H}$ and the zero is the empty relation $0$. For $H=[h]$ define
\begin{equation}
\label{eq:liveness-language}
 \OWL_h=\{R_1\cdots R_\ell\in(\R{[h]})^*:
                  R_1\cdots R_\ell\ne0\}.
\end{equation}
The product of the empty word is $\Id{[h]}$, so $\eps\in\OWL_h$.

\begin{theorem}
\label{thm:main}
For every integer $h\ge2$, the language $\OWL_h$ is recognized by a
nondeterministic automaton with $h+3$ states using no left moves. Every
equivalent deterministic two-way automaton with $s$ states satisfies
\begin{equation}
\label{eq:main-bound}
 4(s+2)^2\ge 2^{\lfloor(h-2)/31\rfloor}.
\end{equation}
If an initially accepting configuration counts as acceptance, $s+2$ in
\eqref{eq:main-bound} may be replaced by $s+1$.
\end{theorem}

\begin{corollary}
\label{cor:no-polynomial}
There are no absolute constants $C,c>0$ such that every $n$-state two-way
nondeterministic automaton over every finite alphabet has an equivalent
two-way deterministic automaton with at most $Cn^c$ states.
\end{corollary}

\begin{proof}
Put $n=h+3$ in Theorem~\ref{thm:main}. Every deterministic equivalent
for this family has at least
\[
 2^{\frac12\lfloor(n-5)/31\rfloor-1}-2
\]
states. This bound exceeds $Cn^c$ for every fixed $C,c$ when $n$ is large
enough. For each $h$ the alphabet has $2^{h^2}$ symbols and is finite.
Since the constants in the proposed simulation must be independent of
the alphabet, the family is within its scope.
\end{proof}

\subsection{Proof architecture}

Write $\D{k}$ for the matching monoid of degree $k$. The proof compares
two bounds on $k$. Theorem~\ref{thm:machine-diagrams} represents each
$s$-state deterministic machine by symbol diagrams in $\D{k}$, with
$k\le4(s+2)^2$. For a machine recognizing $\OWL_h$,
Lemma~\ref{lem:syntactic-quotient} gives a unital surjection from the
submonoid generated by its ordinary-symbol diagrams onto $\R{[h]}$.
Concretely, the map sends each word diagram to its relation product,
preserves multiplication and the empty product, and reaches every
relation on $[h]$. These three properties, including the full image,
are the algebraic input to the lower bound.
Corollary~\ref{cor:divisor-degree} then forces
$k\ge2^{\lfloor(h-2)/31\rfloor}$. These three statements give the deterministic
lower bound in Theorem~\ref{thm:main}.

We prove the algebraic bound first. The \emph{rank} of a matching diagram
counts its pairs joining opposite sides. A diagram $e$ is
\emph{idempotent} if $e^2=e$; its \emph{corner} consists of the diagrams
$z$ with $ez=ze=z$, with identity $e$. Section~\ref{sec:matching} reduces
the corner of an idempotent of rank $r$ to degree $r$, preserving rank
and controlling the indices at which the diagram differs from the
identity. An idempotent of rank $r$ in $\D{k}$ differs
from the identity on at most $2(k-r)$ indices. The support control is
uniform: a family supported in one set remains supported in a single
set of no greater size after reduction. Section~\ref{sec:rank} uses
$32$ conjugates of one idempotent to make $256$ successive relation
additions. A common support set bounds their total rank loss, while
each addition yields a further corner representing all relations on
$h-31$ points. Comparing these bounds doubles the required rank loss
every $31$ points.

Section~\ref{sec:machines} obtains local symbol diagrams by converting
acceptance to reachability of a designated sink. Its undirected component
is a tree, whose doubled tour supplies the matching test. This traversal
has precedents in Sipser's halting method and in the reversible
simulation of Lange, McKenzie, and
Tapp~\cite{SipserHalting1980,LangeMcKenzieTapp2000}. Candidate
transition slots completed by separate leaves make the construction
independent of neighboring symbols. Finally,
Section~\ref{sec:separation} obtains the quotient by using
singleton-identity context letters to test each pair in a relation
product. No input-length or running-time bound is imposed.

\subsection{Related results}

Adeogun and Kapoutsis also analyze limitations of their earlier
lower-bound construction~\cite{AdeogunKapoutsisLimitation2026}. They
prove a quadratic ceiling for its sequences of connectivity properties
and state a general-property extension whose proof is deferred.
The amplification here uses a fixed chain of $256$ additions. Its
exponential bound comes from the inductive lower bound on each rank
loss in a smaller full relation corner, together with one common
support budget for all additions.

The companion paper~\cite[Theorem~1.1]{CompanionComplementation2026}
proves that a nondeterministic two-way automaton for the complement of
$\OWL_h$ requires at least
$\frac12 2^{\lfloor(h-2)/127\rfloor}-1$ states. The present proof
establishes the deterministic bound directly using matching diagrams.
Both results concern state succinctness over
growing alphabets; neither yields a uniform logarithmic-space
separation or a bound measured by transition-table size.
